\scanpage{10}
\mspage{- 85 -}
n\textsuperscript{o} 239

contenu dans $U_\mu$ et qui est telle que la somme $s_\mu+\handwritten{\sum}_{\handwritten{\lambda\in}J}s_{\handwritten{\lambda}}$ soit égale à $1_x$ en tout point $x\handwritten{\in} V_J\handwritten{\cup} V_\mu$ ; on obtiendra ainsi une contradiction, ce qui prouvera le théorème. Posons $s=\handwritten{\sum}_{\handwritten{\lambda\in}J}s_{\handwritten{\lambda}}$ ; les ensembles fermés $\overset{\handwritten{\rule{.45em}{.1ex}}}{V}_\mu$ et $\handwritten{\complement}U_\mu$ étant sans point commun, le th.2 montre qu'il existe une section $t$ de $\handwritten{\unsure{\mathscr A|\mathscr R}}$ au-dessus de $X$ telle que $t(x)=1_x$ pour $x\handwritten{\in}\overset{\handwritten{\rule{.45em}{.1ex}}}{V}_\mu$ et $t(x)=0_x$ pour $x\handwritten{\in}\handwritten{\complement}U_\mu$. Posons $s_\mu(x)=t(x)(1_x-s(x))$ pour tout $x\handwritten{\in} X$; $s_\mu$ est une section au-dessus de $X$ dont le support est contenu dans celui de $t$, donc dans $U_\mu$. D'autre part, on a $s_\mu+s=t(1-s)+s$ ; en tout point de $V_J$ cette somme se réduit à $s(x)=1_x$, et en tout point de $V_\mu$ on a $t(x)=1_x$, donc la somme se réduit à $1_x-s(x)+s(x)=1_x$, ce qui achève la démonstration.

6. \uline{Faisceaux fins}.

DEFINITION 2. - \uline{Etant donné un espace localement paracompact} $X$, \uline{on dit qu'un faisceau de} $\handwritten{\mathscr A}$-\uline{modules} $\handwritten{\mathscr F}$ \uline{de base} $X$ \uline{est fin si le faisceau d'anneaux} $\operatorname{Hom}_{\handwritten{\mathscr A}}(\handwritten{\mathscr F},\handwritten{\mathscr F})$ (\S\,2,n\textsuperscript{o}7) \uline{est mou}.

En vertu de la prop.3, et de la définition des sections dans $\operatorname{Hom}(\handwritten{\mathscr F},\handwritten{\mathscr F})$ (\S\,2,n\textsuperscript{o}7), il revient au même de dire qu'il existe une base $\handwritten{\unsure{\phi|\Phi}}$ de la topologie de $X$ formée d'ensembles relativement paracompacts et, pour chaque $U\handwritten{\in\phi}$, un ensemble $\handwritten{\psi}(U)$ de parties fermées de $X$ contenues dans $U$, vérifiant $(\mathrm{FM}_{\mathrm I})$ et la condition

$(\mathrm{FM}^{\prime\prime}_{\mathrm{II}})$ \uline{Quels que soient} $U\handwritten{\in\phi}$, $F\handwritten{\in\psi}(U)$ \uline{et le voisinage ouvert} $V$ \uline{de} $F$, \uline{contenu dans} $U$, \uline{il existe un endomorphisme} $\theta$ \uline{de} $\handwritten{\mathscr F}$ \uline{qui soit l'identité au-dessus d'un voisinage convenable de} $F$ \uline{et se réduise à} $0$ \uline{au-dessus d'un voisinage convenable de} $\handwritten{\complement}V$.

PROPOSITION 4. - \uline{Tout faisceau fin de} $\handwritten{\mathscr A}$-\uline{modules} $\handwritten{\mathscr F}$ \uline{est un faisceau mou}.
% archive sheet number: 9

